\documentclass[11pt,a4paper]{article}
  
\usepackage{amsmath}
\usepackage{amssymb}
\usepackage{amsfonts}
\usepackage[dvips]{graphicx}
\usepackage{color}

% \usepackage[catalan]{babel}
\usepackage[latin1]{inputenc}

\definecolor{titol}{rgb}{0.2, 0.7, 0.8}



\pagestyle{empty}
\setlength{\textwidth}{16.5 cm}
\addtolength{\oddsidemargin}{-2.5 cm}
\addtolength{\topmargin}{-3 cm}
\setlength{\textheight}{26 cm}
\nofiles

\newcounter{probl}
\newcounter{cs}
\stepcounter{cs}


\newenvironment{pr}{\stepcounter{probl}\noindent{\llap{\bf\theprobl .\hspace{.1cm}}}}%
                   {\vspace{.5cm}}
\newcommand{\Sol}{\noindent \underline{\bf Soluci\'{o} al problema \theprobl}:}         
%%% Carles: Afegit \Sol, mogut "\stepcounter{probl}" al savant de l'entorn "pr" de manera que 
%%%           podem comen�ar sense donor el valor 1 probl.        

\newcommand{\p}{\begin{pr}}
\newcommand{\fp}{\end{pr}}

\newcommand{\casos}{\begin{itemize}}
\newcommand{\fcasos}{\end{itemize}\setcounter{cs}{1}}
\newcommand{\cas}{\item[(\alph{cs})]\stepcounter{cs}}

\newif\ifsolucions
%% Descomentar aix� per veure solucions
\solucionstrue

\newcommand{\newpagewithsolutions}{\ifsolucions\newpage\fi} %%% Added Carles 18/09/2012

\newcommand{\C}{\mathbb{C}}
\newcommand{\N}{\mathbb{N}}
\newcommand{\R}{\mathbb{R}}
\newcommand{\Q}{\mathbb{Q}}
\newcommand{\T}{\mathbb{T}}
\newcommand{\Z}{\mathbb{Z}}




%
%
%
%\documentclass[11pt, oneside]{article}   	% use "amsart" instead of "article" for AMSLaTeX format
%\usepackage{geometry}                		% See geometry.pdf to learn the layout options. There are lots.
%%\geometry{A4}  
%\usepackage[latin1]{inputenc}
%\usepackage[catalan]{babel}
%
%
%
%
%
%
%                 		% ... or a4paper or a5paper or ... 
%%\geometry{landscape}                		% Activate for rotated page geometry
%%\usepackage[parfill]{parskip}    		% Activate to begin paragraphs with an empty line rather than an indent
%\usepackage{graphicx}				% Use pdf, png, jpg, or eps§ with pdflatex; use eps in DVI mode
%								% TeX will automatically convert eps --> pdf in pdflatex		
%\usepackage{amssymb}
%
% % % \usepackage{termcal}
% % \makeatletter
% % %%%
% % \def\setdate@#1/#2/#3!{
% % \setcounter{month}{#2}
% % \setcounter{date}{#1}
% % \setcounter{year}{#3}
% % \global\newmonthtrue\setleap}
% % %%%
% % \renewcommand{\monthname}{\ifcase\c@month\or Gen\or Feb\or Mar�\or Abr%
% % \or Maig\or Juny\or Jul\or Ago\or Set\or Oct%
% % \or Nov\or Des\fi}
% % %%%
% % \renewcommand{\ordinaldate}{\arabic{date}}  %?????????????????????
% % %\renewcommand\ordinal[1]{% 
% % %     \let\last@=\relax\let\last@@=\relax 
% % %      \expandafter\@rd\the#1x}
% % %\renewcommand\@rd[1]{\ifx#1x\if\last@@1th\else\@rdend{\last@}\fi\else 
% % %       \let\last@@=\last@\def\last@{#1}#1\expandafter\@rd\fi}
% % %       
% % %\renewcommand\@rdend[1]{\ifcase#1 th\or st\or nd\or rd\else th\fi}
% % %%%
% % \makeatother
% % %%%%%%%%%%%%%%%%%%
% % \renewcommand\curdate{\arabic{date}/\arabic{month}/\arabic{year}}
% % \renewcommand{\calprintdate}{%
% %     \ifnewmonth\framebox{\ordinaldate\ de \monthname}%
% %      \else \ordinaldate\fi
% % }
% % %%%%%%%%%%%%%%%%%%%
% % 
% % 
% % \newcommand{\TRClass}{%
% % \skipday % Tuesday
% % \skipday % Tuesday
% % \skipday % Wednesday (no class)
% % \calday[Thursday]{\classday} % Thursday
% % \skipday % Friday 
% % \skipday\skipday % weekend (no class)
% % }
% % %
% % \newcommand{\Holiday}[2]{%
% % \options{#1}{\noclassday}
% % \caltext{#1}{#2}
% % }
% % %SetFonts

%SetFonts


\begin{document}

%%%%%%%% cap�alera %%%%%%%%%%%%%%%%%%%%%%%%%%%%%
\hrule\smallskip \medskip \noindent\textbf{Estad\'{i}stica
\hfill\textbf{Enginyeria Inform�tica}}\par\par
                    \medskip\noindent\textsl{\bf Curs 2026 -- 2027}
\hfill\textbf{UAB}
                     \medskip\hrule\vspace{0.5cm} % CAP\c{C}ALERA
%\vspace{0.2cm}
\vspace{0.5cm}
%%%%%%%%%%%%%%%%%%%%%%%%%%%%%%%%%%%%%%%%%%%%%%%


\centerline{\bf Proposta de programaci� (�s intern, no es mostra als estudiants)}

\bigskip

\hspace{-2cm}
\begin{tabular}{|c|l|c|l|}
\hline
&Dimarts (problemes) & & Divendres (teoria)
\\
\hline
\hline

 15 Set &
  
%  \begin{tabular}{l}  %
    Presentaci\'o de l'assignatura\\
    i tema 0
%  (V�deo)
%  \end{tabular}
 & 
18 Set &
 \begin{tabular}{l}  
   \bf Descriptiva 1\\ \small Poblaci\'o, mostra, variables,  %\\ \small
   Taules de Freq\"u\`encia, 
 gr\`afiques
 \\ \small Estudi descriptiu d'una variable. 
\\ \small 
Mesures de centre i dispersi\'o \\
 \end{tabular}
\\
\hline
\hline

 22 Set &
  
 Problemes descriptiva 1 (llista 1)
%  \begin{tabular}{l}  %
%     Presentaci\'o de l'assignatura %\\
%  (V�deo)
%  \end{tabular}
 & 
25 Set &
  \begin{tabular}{l}   
   \bf Descriptiva 2\\
\small Estudi descriptiu de dues variables.\\
\small Taules de conting\`encia.
\small Regressi\'o lineal 
\end{tabular}
\\
\hline
29 Set &
Problemes descriptiva 2 (llista 1)
& 2 Oct &
% \begin{tabular}{l}   
%    \bf Descriptiva 2\\
% \small Estudi descriptiu de dues variables.\\
% \small Taules de conting\`encia.
% \small Regressi\'o lineal 
% \end{tabular}
\begin{tabular}{l}\bf Probabilitat 1\\ \small Introducci\'o, Taules de freq\"u\`encia i probabilitat \\
\small Probabilitat condicionada, Independ\`encia d'esdeveniments \\ \small i F\'ormula de Bayes.
\end{tabular}
\\
\hline
6 oct &    Problemes probabilitat 1 (llista 2)
 &
 9 oct &
\begin{tabular}{l}\bf Probabilitat 2\\ \small Variables aleat\`ories, esperan\c ca i vari\`ancia \\
\small Distribucions discretes: Bernoulli, Binomial i Poisson.
\end{tabular}
\\
\hline
13 Oct &
 \begin{tabular}{l}
 Problemes probabilitat 2 (llista 2)
          \\ (1a entrega de problemes)
         \end{tabular}
          & 16 Oct &

\begin{tabular}{l}\bf Probabilitat 3\\ \small Distribucions cont\'\i nues: Uniforme, exponencial i Normal  
% % \\
% % \small Aproximaci\'o de la binomial per la normal.
\end{tabular}
\\
\hline
20 Oct & Problemes variables aleat�ries 1 (llista 3)   &
23 Oct &
\begin{tabular}{l}\bf Probabilitat 4 \\ \small Teorema Central del L\'\i mit \\
 \small Aproximaci\'o de la binomial per la normal.
 \\ \small Fer algun problema de la llista 3?
\end{tabular}
\\
\hline
27 Oct &        &
30 Oct &
\\
\hline
3 Nov &
  & 6 Nov &
\begin{tabular}{l}\bf Infer\`encia 1\\ \small
\small Intervals de Confian\c ca I: Definici\'o, \\
\small  Intervals per la mitjana, t-Student
\end{tabular}
  \\
\hline
10 Nov & Problemes variables aleat�ries 2 (llista 3)  &
13 Nov &
\begin{tabular}{l}\bf   Infer\`encia 2\\ 
\small Intervals de Confian\c ca II: Intervals per la vari\`ancia. \\
\small   Intervals per les proporcions 
\end{tabular}
% Examen Parcial
\\
\hline
17 Nov & Problemes d'intervals de confian�a 1 (llista 4) &
20 nov &
 \begin{tabular}{l}\bf Infer\`encia 3\\ 
\small Test d'hip\`otesi I: concepte de test d'hip\`otesi \\
\small  test per a l'esperan\c ca, vari\`ancia i proporci\'o 
\end{tabular}
\\
\hline
24 Nov &
 Problemes IC 2 (llista 4)
& 27 nov &
\begin{tabular}{l}\bf Infer\`encia 4\\ 
\small Test d'hip\`otesi II: comparaci\'o de dues poblacions \\
\small  Dades aparellades, Mostres independents. \\
\small Test de comparaci\'o de mitjanes
\end{tabular}
\\
\hline
1 Des & \begin{tabular}{l} Problemes IC 3 (llista 4)
 \\
 (2a entrega de problemes)
\end{tabular} &
4 Des
& 
\begin{tabular}{l}\bf Infer\`encia 5\\
\small Test d'hip\`otesi III: test de comparaci\'o de vari\`ancies \\
\small tests no param\`etrics: taules de conting\`encia, \\  \small Test de bondat d'ajust \\
\end{tabular}
\\
\hline
8 Des &
& 11 Des &
\begin{tabular}{l}\bf Infer\`encia 6\\
\small Rep�s/completar  \\
\end{tabular}
\\
\hline
15 Des & \begin{tabular}{l}
 Problemes Test d'Hip�tesi 1 (llista 5)
             \\ 3a entrega de problemes
         \end{tabular}
 & 18 Des   &  \begin{tabular}{l}
 Problemes Test d'Hip�tesi  2 i 3 (llista 5)
             \\ 3a entrega de problemes
         \end{tabular}
\\
%
\hline
22 Des &
 &    &
\\
%
\hline
\end{tabular}

\begin{itemize}
 \item Primer parcial:    / Examen final:
 /
 Recuperaci\'o:
\end{itemize}





\end{document}
